\section{مقدمه و هدف}
\label{sec:intro}
همان‌طور که در بخش قبلی گزارش دیدید، شناساگر (\lr{tagger}) جت‌های بوست‌شده \Zbb را طراحی کرده‌ایم. هدف این بخش تولید داده مناسب برای تشخیص رویدادهای \ttZ از پس‌زمینه \ttbb است؛ به این صورت که شناساگر جت‌های \lr{Z} را مشخص می‌کند و جداکننده رویداد (\lr{event discriminator}) با خروجی آن رویدادهای \ttZ را از \ttbb جدا می‌کند. برای اینکه تفاوت موجود در توزیع‌ها حاصل از روش تولید متفاوت نباشد و تنها ناشی از فیزیک فرایندها باشد، همه رویدادها را در یک چارچوب شبیه‌سازی واحد تولید کردیم: نسخه نرم‌افزارها، تعریف جت، برش‌های تحلیل، برچسب منشأ جت‌ها و قالب ذخیره‌سازی برای همه نمونه‌ها یکی است. از آنجا که شناساگر روی یک داده آموزش می‌بیند ولی خروجی آن را روی داده دوم، یعنی جت‌های رویدادهای \ttZ و \ttbb، می‌خواهیم، تمام تلاشمان این بود که پردازش دو داده یکسان باشد تا خروجی روی داده دوم تا حد ممکن منصفانه باشد. با این حال، این دو داده فرایندهای فیزیکی متفاوتی هستند، بنابراین توزیع‌های آنها همچنان کمی با هم فرق دارند.

برای ورودی شناساگر \lr{Z} یک نمونه سیگنال می‌سازیم که \lr{Z+jet} است: یک بوزون \lr{Z} پرتکانه همراه با یک جت تولید می‌شود و \lr{Z} به یک زوج \bb واپاشی می‌کند که در یک جت بوست‌شده جمع می‌شود. نمونه پس‌زمینه \lr{bbbb} است: چهار کوارک \lr{b} (دو زوج \bb) در برهمکنش \lr{QCD} تولید می‌شوند و جت بوست‌شده از یکی از این زوج‌ها ساخته می‌شود، یعنی جتی که \bb دارد ولی از \lr{Z} نیامده است. وقتی شناساگر \lr{Z} بین این دو نمونه تمایز ایجاد کند، خروجی آن به یک جداکننده کوانتومی داده می‌شود که سیگنال \ttZ را، یعنی تولید یک زوج کوارک تاپ همراه با بوزون \lr{Z} که به \bb وامی‌پاشد و علاوه بر جت \lr{Z} جت‌های ناشی از واپاشی تاپ‌ها را دارد، از پس‌زمینه \ttbb جدا کند؛ در \ttbb یک زوج تاپ همراه با یک زوج \bb از برهمکنش \lr{QCD} تولید می‌شود. از هر نمونه ۲۰۰۰۰ رویداد در برخورد \lr{pp} با انرژی \lr{13~TeV} ساخته شده است.

پرسش اصلی این است که آیا متغیرهای زیرساختار جت و ویژگی‌های ذره‌ها برای هر چهار نمونه رفتار هماهنگی دارند، و آیا شناساگری که روی یک جفت (سیگنال، پس‌زمینه) آموزش دیده روی جفت دیگر هم کار می‌کند.

در ادامه، بخش~\ref{sec:sim} زنجیره شبیه‌سازی، برش‌ها و برچسب منشأ جت‌ها را توضیح می‌دهد؛ بخش~\ref{sec:vars} متغیرهای جت و ذره را تعریف می‌کند. در بخش~\ref{sec:cutflow} تعداد جت‌ها بعد از هر برش آمده است و در بخش‌های~\ref{sec:set1} تا~\ref{sec:particles} نمودارهای چهار نمونه را کنار هم می‌گذاریم و مقایسه می‌کنیم. بخش~\ref{sec:disc} به بحث و محدودیت‌ها و بخش~\ref{sec:next} به گام بعدی اختصاص دارد.
